一个模型上线后的头两周，报障单常常是这几条：拦截量是测算的四倍；阈值一直用着默认的 \(0.5\)；模型说 \(0.9\) 的那批人实际只有六成出事；人工复核队列每天积压七十条。四条里没有一条与模型的拟合能力有关。它们全部发生在概率算出来之后。

这一单元处理的就是那之后的事。把问题写成决策的形状需要三样东西：未知状态、可选行动、损失函数。有了它们，风险是损失的期望，规则之间可以比较，阈值有闭式解。缺了损失函数不等于没有损失，只等于默认了两类错误一样贵——分类里那条 \(0.5\) 的判定线就是这么来的，而漏检与误报的实际代价常常差着两三个数量级。

概率本身也要被检验。评分规则决定模型有没有动机把信念如实报出来，Brier 与对数损失都让诚实报告成为期望最优；校准则是另一回事，说 \(70\%\) 的那些事件里应当约有 \(70\%\) 发生。区分度好的模型可以严重失准，完美校准的常数模型可以毫无用处，两者互不蕴含，得分开报。再往前一步，点概率回答不了``这一批结论里有多少会落空''，conformal prediction 用一条可交换性假设换来有限样本、无分布假设的覆盖保证，它承诺的是边缘覆盖，不承诺每个子群、每个个体都达标。

最后一步是把这些接到现实的处置能力上。复核团队每天只能看两百条时，问题从``该不该拦这一条''变成``这一批里选哪两百条''，前面的损失最小化加一个预算约束就得到答案，排序的对象是单位成本的净收益。

六节依次展开：

\begin{enumerate}
\tightlist
\item
  状态、行动、损失与风险构成一个决策问题——用同一套语言重写估计、检验与预测，说明风险是一整条曲线，以及损失函数不写下来时系统会替你填一个。
\item
  Bayes rule 最小化后验期望损失——把``给出概率''与``做出决定''拆成两步，给出由损失比决定的阈值公式，并算几个远离 \(0.5\) 的具体例子。
\item
  Minimax 与可容许性比较整条风险曲线——最坏情形与加权平均读的是同一条曲线的不同部位，可容许性只是一道最低门槛，两种准则各自对应什么样的部署场景。
\item
  Proper scoring rule 与校准衡量概率是否诚实——从 Brier 分解得到严格 proper 的含义，把区分度与校准分开，并给出可靠性图的画法与读法。
\item
  Conformal prediction 用可交换性换有限样本覆盖——秩的均匀性如何直接给出覆盖不等式，它保证什么、不保证什么，以及漂移、时序与数据复用怎样让保证失效。
\item
  拒绝预测与容量约束把概率变成行动——拒绝带的两条边界来自复核成本，容量约束把逐点阈值变成批内排序，被截断的净收益就是扩容的影子价格。
\end{enumerate}

这一单元不假定存在一个脱离任务的``最可信模型''。同一个概率模型在对数损失下表现优秀，换到罕见灾难的非对称损失上可能完全不能用；同一个 conformal 集合达到了名义覆盖，也可能宽到没有任何决策价值。谈可信时必须说清对象：哪一种随机机制、哪一个目标人群、哪一类行动、哪一种错误。

\subsection{怎么读这一章}

前三篇建立框架，后三篇转向概率预测的可信性，最后一篇把两条线合起来。手上正有模型待上线的读者可以直接从第二篇的阈值公式和第四篇的可靠性图入手，需要交付区间或集合时再回到第五篇。读每一篇时反复追问同一件事：这个数字对应的是哪个人群，代价由谁承担，以及行动做出之后谁来兑现。

\subsection{本章篇目}

以下篇目按概念依赖排列；若从具体问题进入，也可以先打开最接近当前困惑的一篇，再沿篇内链接回补。

\begin{itemize}
\tightlist
\item \hyperref[sec:11-Mathematical-Statistics/07-Statistical-Decision-and-Reliable-Prediction/01]{``状态、行动、损失与风险构成一个决策问题''}；
\item \hyperref[sec:11-Mathematical-Statistics/07-Statistical-Decision-and-Reliable-Prediction/02]{``Bayes rule 最小化后验期望损失''}；
\item \hyperref[sec:11-Mathematical-Statistics/07-Statistical-Decision-and-Reliable-Prediction/03]{``Minimax 与可容许性比较整条风险曲线''}；
\item \hyperref[sec:11-Mathematical-Statistics/07-Statistical-Decision-and-Reliable-Prediction/04]{``Proper scoring rule 与校准衡量概率是否诚实''}；
\item \hyperref[sec:11-Mathematical-Statistics/07-Statistical-Decision-and-Reliable-Prediction/05]{``Conformal prediction 用可交换性换有限样本覆盖''}；
\item \hyperref[sec:11-Mathematical-Statistics/07-Statistical-Decision-and-Reliable-Prediction/06]{``拒绝预测与容量约束把概率变成行动''}。
\end{itemize}
